\section*{第 \ref{ch:b1:respiration-fermentation} 章——细胞呼吸与发酵}

\begin{solution}{exo:b1:respiration-fermentation:1}
\omterm{def:b1:respiration-fermentation:glycolysis}{糖酵解}，\omterm{def:b1:eukaryotic-cell:organelle}{细胞溶胶}：2 分子丙酮酸、2 个 \omterm{def:b1:water-small-molecules:atp}{ATP}、2 个 NADH。丙酮酸氧化与
\omterm{def:b1:respiration-fermentation:krebs}{克雷布斯循环}，\omterm{def:b1:eukaryotic-cell:mitochondrion}{线粒体}基质：6 个 $\mathrm{CO_2}$、8 个 NADH、2
个 $\mathrm{FADH_2}$、2 个 GTP。氧化磷酸化，内膜：
由氧生成水，约 26 个 \omterm{def:b1:water-small-molecules:atp}{ATP}。
\end{solution}

\begin{solution}{exo:b1:respiration-fermentation:2}
葡萄糖 $+ 2\,\mathrm{NAD^+} + 2\,\mathrm{ADP} + 2\,\mathrm{P_i} \to
2$ 丙酮酸 $+ 2\,\mathrm{NADH} + 2\,\mathrm{H^+} + 2\,\mathrm{ATP} +
2\,\mathrm{H_2O}$。\omterm{def:b1:respiration-fermentation:glycolysis}{底物水平磷酸化}：一个磷酸基团直接
由高能中间产物（1,3-二磷酸甘油酸、磷酸烯醇式丙酮酸）
经酶转移给 ADP，
不依赖膜梯度，也不需要氧。
\end{solution}

\begin{solution}{exo:b1:respiration-fermentation:3}
NADH 与 $\mathrm{CO_2}$（异柠檬酸至 $\alpha$-酮戊二酸），NADH
与 $\mathrm{CO_2}$（$\alpha$-酮戊二酸至琥珀酰\omterm{def:b1:enzymes:enzyme}{辅酶} A），GTP
（琥珀酰\omterm{def:b1:enzymes:enzyme}{辅酶} A 至琥珀酸），$\mathrm{FADH_2}$（琥珀酸至延胡索酸），
NADH（苹果酸至草酰乙酸）。两个 $\mathrm{CO_2}$ 来自
于六碳酸与五碳酸的两次脱羧。
\end{solution}

\begin{solution}{exo:b1:respiration-fermentation:4}
它使 $\mathrm{NAD^+}$ 再生，以使\omterm{def:b1:respiration-fermentation:glycolysis}{糖酵解}得以继续，
并维持每葡萄糖产生 2 个 \omterm{def:b1:water-small-molecules:atp}{ATP}。它不氧化碳，
不提取剩余能量（产物几乎保留全部能量），也
不能产生超过呼吸作用十五分之一的 \omterm{def:b1:water-small-molecules:atp}{ATP}。
\end{solution}

\begin{solution}{exo:b1:respiration-fermentation:5}
至氧：$-2\times 96\,500\times 1.14 = \qty{-220}{kJ/mol}$，四个 \omterm{def:b1:water-small-molecules:atp}{ATP}（
每个 \qty{50}{kJ}）。至泛醌：$-2\times 96\,500\times 0.37 =
\qty{-71}{kJ/mol}$，一个 \omterm{def:b1:water-small-molecules:atp}{ATP}。
\end{solution}

\begin{solution}{exo:b1:respiration-fermentation:6}
肌肉：$2 + 2$（GTP）$+ 8\times 2.5$（\omterm{def:b1:eukaryotic-cell:mitochondrion}{线粒体} NADH）$+ 2\times
1.5$（$\mathrm{FADH_2}$）$+ 2\times 1.5$（经 \omterm{def:b1:respiration-fermentation:carriers}{FAD} 转移的\omterm{def:b1:eukaryotic-cell:organelle}{细胞溶胶} NADH）$=
30$。肝脏：\omterm{def:b1:eukaryotic-cell:organelle}{细胞溶胶} NADH 给出 $2\times 2.5$：32。
\end{solution}

\begin{solution}{exo:b1:respiration-fermentation:7}
RQ $= 2.2$。经呼吸消耗的葡萄糖 $x$：$6x = 1.0$，$x = \qty{0.167}{mmol}$；
$\mathrm{CO_2}$：$6x + 2y = 2.2$，$2y = 1.2$，$y = \qty{0.6}{mmol}$。
\omterm{def:b1:respiration-fermentation:fermentation}{发酵}的葡萄糖占 $0.6/0.767 = \qty{78}{\%}$，
呼吸的占 \qty{22}{\%}。
\end{solution}

\begin{solution}{exo:b1:respiration-fermentation:8}
耗氧量上升（\omterm{def:b1:respiration-fermentation:chain}{呼吸链}摆脱了梯度的反压，
全速运行），\omterm{def:b1:water-small-molecules:atp}{ATP} 合成停止（无梯度驱动
合酶），电子的全部能量都转化为
热量。患者迅速燃烧\omterm{def:b1:lipids:triglyceride}{脂肪}，死于
无法降下的体温。
\end{solution}

\begin{solution}{exo:b1:respiration-fermentation:9}
\omterm{def:b1:respiration-fermentation:chain}{呼吸链}发生堵塞：各载体均被还原，没有质子
被泵出，不产生 \omterm{def:b1:water-small-molecules:atp}{ATP}。NADH 不再被氧化，因此\omterm{def:b1:respiration-fermentation:krebs}{克雷布斯循环}
与丙酮酸脱氢酶因缺乏 $\mathrm{NAD^+}$ 而停止。
\omterm{def:b1:respiration-fermentation:glycolysis}{糖酵解}解除了 \omterm{def:b1:water-small-molecules:atp}{ATP} 的抑制，快速进行，必须
再生其 $\mathrm{NAD^+}$，方法是将丙酮酸还原为乳酸，乳酸
大量涌入血液。
\end{solution}

\begin{solution}{exo:b1:respiration-fermentation:10}
$8\times 10 + 7\times 2.5 + 7\times 1.5 - 2 = 80 + 17.5 + 10.5 - 2 =
106$ 个 \omterm{def:b1:water-small-molecules:atp}{ATP}；每个碳 $106/16 = 6.6$ 个。氧耗：棕榈酸 23
个 $\mathrm{O_2}$ 对应 106 个 \omterm{def:b1:water-small-molecules:atp}{ATP}，即 \qty{0.22}{O_2} 每个 \omterm{def:b1:water-small-molecules:atp}{ATP}；葡萄糖 6 个对应 30
个 \omterm{def:b1:water-small-molecules:atp}{ATP}，即 \qty{0.20}{O_2} 每个 \omterm{def:b1:water-small-molecules:atp}{ATP}。以单位氧计，葡萄糖
高出 \qty{10}{\%}，水下缺氧的海豹偏好碳水化合物；以单位质量计，
\omterm{def:b1:lipids:triglyceride}{脂肪}的效能约为两倍：飞行中受质量限制的蜂鸟，
迁徙靠\omterm{def:b1:lipids:triglyceride}{脂肪}供能。
\end{solution}

\begin{solution}{exo:b1:respiration-fermentation:11}
仅 \omterm{def:b1:water-small-molecules:atp}{ATP}：$5/3 < \qty{2}{s}$。磷酸肌酸：$25/3 = \qty{8}{s}$。
\omterm{def:b1:respiration-fermentation:fermentation}{发酵}：每秒 \omterm{def:b1:water-small-molecules:atp}{ATP} 生成量为 $2\times 3 = \qty{6}{mmol/L}$，足以
满足需求，只要\omterm{def:b1:carbohydrates:polysaccharide}{糖原}与对乳酸的耐受能够维持
（数十秒）。氧的供应需要一分钟或更久才能上升，
而在肌纤维中每秒至多提供约 \qty{1}{mmol/L} 的 \omterm{def:b1:water-small-molecules:atp}{ATP}：
十秒的短跑在呼吸作用开始发挥作用之前已经结束，并
由磷酸肌酸与\omterm{def:b1:respiration-fermentation:fermentation}{发酵}供能。
\end{solution}

\begin{solution}{exo:b1:respiration-fermentation:12}
梯度以跨绝缘膜的质子浓度差与
电势差的形式储存能量——恰似
带电电容器加浓差电池；\omterm{def:b1:respiration-fermentation:chain}{呼吸链}充当充电器，
由电子下降驱动；\omterm{def:b1:water-small-molecules:atp}{ATP} 合酶是一台由
质子流推动的马达，恰似涡轮机，将流动耦合为
机械旋转从而生成 \omterm{def:b1:water-small-molecules:atp}{ATP}。此类比并不严密之处在于“电池”
只能维持\omterm{def:b1:cell-unit-of-life:cell}{细胞}数毫秒的消耗，且须
持续充电；同一膜上的其他转运蛋白也会耗用它，
并且涡轮机可反向运转，水解 \omterm{def:b1:water-small-molecules:atp}{ATP}，
以在\omterm{def:b1:respiration-fermentation:chain}{呼吸链}失效时泵送质子。
\end{solution}

\begin{solution}{pb:b1:respiration-fermentation:1}
\textbf{1.} \omterm{def:b1:respiration-fermentation:glycolysis}{糖酵解}：2 个 \omterm{def:b1:water-small-molecules:atp}{ATP}、2 个 NADH（\omterm{def:b1:eukaryotic-cell:organelle}{细胞溶胶}中）。丙酮酸
脱氢酶：2 个 NADH。\omterm{def:b1:respiration-fermentation:krebs}{克雷布斯循环}：6 个 NADH、2 个 $\mathrm{FADH_2}$、2 个 GTP。
\textbf{2.} \omterm{def:b1:eukaryotic-cell:mitochondrion}{线粒体} NADH $8\times 10 = 80$；$\mathrm{FADH_2}$
$2\times 6 = 12$；\omterm{def:b1:eukaryotic-cell:organelle}{细胞溶胶} NADH 按 $\mathrm{FADH_2}$ 计 $2\times 6 = 12$：
共 104 个质子。
\textbf{3.} $104/4 = 26$ 个 \omterm{def:b1:water-small-molecules:atp}{ATP}；总计 $26 + 2 + 2 = 30$。
\textbf{4.} 2 个 \omterm{def:b1:water-small-molecules:atp}{ATP}。
\textbf{5.} 15.
\textbf{6.} 有氧时 $30\times 50 = \qty{1500}{kJ}$，占
2870 的 \qty{52}{\%}。无氧时：\qty{100}{kJ}，即 2870 的 \qty{3.5}{\%}，但
却占所释放的 235 中的 \qty{43}{\%}——\omterm{def:b1:respiration-fermentation:fermentation}{发酵}善于捕获其所释放的能量，
只是释放得很少。
\textbf{7.} 通气时每小时消耗葡萄糖 $30/30 = \qty{1}{mmol}$；密封时
$30/2 = \qty{15}{mmol}$.
\textbf{8.} 密封时每小时产生 \qty{30}{mmol} 乙醇与 \qty{30}{mmol} 的
$\mathrm{CO_2}$。通气时：\qty{6}{mmol} 的 $\mathrm{CO_2}$
被释放，并消耗 \qty{6}{mmol} 的 $\mathrm{O_2}$。
\textbf{9.} 每摩尔 $2\times 10.5 = \qty{21}{g}$，
即每克葡萄糖 \qty{0.12}{g}。
\textbf{10.} $0.5\times 180 = \qty{90}{g}$ 每摩尔，即 $90/10.5 =
8.6$ 个 \omterm{def:b1:water-small-molecules:atp}{ATP} 的量，对比可用的 30：\omterm{def:b1:water-small-molecules:atp}{ATP} 过剩，产量
受碳限制（一半葡萄糖被氧化为 $\mathrm{CO_2}$ 以
生成 \omterm{def:b1:water-small-molecules:atp}{ATP}，而生物量需要碳骨架与还原力），
多余的 \omterm{def:b1:water-small-molecules:atp}{ATP} 用于维持消耗。
\textbf{11.} 要获得相同的 \omterm{def:b1:water-small-molecules:atp}{ATP}，密封培养物消耗的葡萄糖
多十五倍（第 7 题），且每克的生长量只有
五分之一（第 9 题）：糖消耗快，生长少，产生酒精——
巴斯德的观察。磷酸果糖激酶受 \omterm{def:b1:water-small-molecules:atp}{ATP} 与
柠檬酸的抑制（通气条件下），呼吸停止时即被解除。
\textbf{12.} $\mathrm{CO_2}$ $44/44 = \qty{1.0}{mmol}$；$\mathrm{O_2}$
$20/32 = \qty{0.625}{mmol}$：RQ $= 1.6$。
\textbf{13.} $\mathrm{O_2}$：$6x = 0.625$，$x = \qty{0.104}{mmol}$。
$\mathrm{CO_2}$：$6x + 2y = 1.0$，$2y = 0.375$，$y = \qty{0.1875}{mmol}$。
\textbf{14.} \omterm{def:b1:respiration-fermentation:fermentation}{发酵}的葡萄糖占 $0.1875/0.292 = \qty{64}{\%}$。\omterm{def:b1:respiration-fermentation:fermentation}{发酵}产生的
\omterm{def:b1:water-small-molecules:atp}{ATP} 为 $2y = 0.375$，呼吸产生的为 $30x = 3.12$，约占
\qty{11}{\%}.
\textbf{15.} $16/23 = 0.70$.
\textbf{16.} 禁食时动物燃烧\omterm{def:b1:lipids:triglyceride}{脂肪}（RQ 0.7）；进食后，它燃烧
这餐中的\omterm{def:b1:carbohydrates:monosaccharide}{糖类}（1.0）。
\textbf{17.} $96\,500\times 0.22 = \qty{21.2}{kJ/mol}$.
\textbf{18.} $4\times 21.2 = \qty{85}{kJ}$，相对每个 \qty{50}{kJ} 的 \omterm{def:b1:water-small-molecules:atp}{ATP}：
\qty{59}{\%}.
\textbf{19.} $2\times 96\,500\times 1.14 = \qty{220}{kJ}$；十个质子
储存 $10\times 21.2 = \qty{212}{kJ}$：\qty{96}{\%}——\omterm{def:b1:respiration-fermentation:chain}{呼吸链}
浪费很少；损失发生在合酶与转运蛋白处。
\textbf{20.} $0.96\times 0.59 = 0.57$；$2.5\times 50/220 = 0.57$：
两种途径得出同一数值。
\textbf{21.} 耗氧量上升，\omterm{def:b1:water-small-molecules:atp}{ATP} 产量下降，趋向
\omterm{def:b1:respiration-fermentation:fermentation}{发酵}的水平（\omterm{def:b1:respiration-fermentation:glycolysis}{糖酵解}每葡萄糖产 2 个 \omterm{def:b1:water-small-molecules:atp}{ATP}，且磷酸果糖激酶被解除抑制，因此
葡萄糖消耗增加），产热上升，生长骤降。
\textbf{22.} \omterm{def:b1:water-small-molecules:atp}{ATP} 合成停止；梯度无法释放，
\omterm{def:b1:respiration-fermentation:chain}{呼吸链}因此受阻，耗氧停止，\omterm{def:b1:cell-unit-of-life:cell}{细胞}则尽量
\omterm{def:b1:respiration-fermentation:fermentation}{发酵}。加入解偶联剂后，质子漏回，\omterm{def:b1:respiration-fermentation:chain}{呼吸链}
与耗氧恢复，但仍无 \omterm{def:b1:water-small-molecules:atp}{ATP} 产生。
\textbf{23.} NADH 仅经复合体 III 和 IV 泵出质子：$6/4 = 1.5$；
每葡萄糖：$8\times 1.5 + 2\times 1.5 + 2\times 1.5 + 4 = 22$ 个 \omterm{def:b1:water-small-molecules:atp}{ATP}。
\textbf{24.} 无氧时\omterm{def:b1:respiration-fermentation:chain}{呼吸链}无法接受电子，因此 NADH
不被再氧化，NADH 池维持还原态；循环中的三个
依赖 \omterm{def:b1:respiration-fermentation:carriers}{NAD} 的脱氢酶没有 $\mathrm{NAD^+}$ 可用，
循环便停止——这正是\omterm{def:b1:respiration-fermentation:fermentation}{发酵}须在\omterm{def:b1:eukaryotic-cell:organelle}{细胞溶胶}中
再生 $\mathrm{NAD^+}$ 以使\omterm{def:b1:respiration-fermentation:glycolysis}{糖酵解}继续进行的原因。
\textbf{25.} 有氧时 30 个 \omterm{def:b1:water-small-molecules:atp}{ATP}，无氧时 2 个，比值为 15；P/O 对 NADH 为 2.5
（10 个质子除以每个 \omterm{def:b1:water-small-molecules:atp}{ATP} 的 4 个），对 $\mathrm{FADH_2}$ 为 1.5（6 个除以
4).
\end{solution}
